\section*{अध्याय \ref{ch:b2:muscle-movement} --- पेशी-संकुचन और प्रेरक कार्य}

\begin{solution}{exo:b2:muscle-movement:1}
पतले तंतु मोटों के साथ सरकते हैं, और उन्हें मायोसिन शीर्ष खींचते हैं;
कोई भी तंतु छोटा नहीं होता। I पट्टी तथा H क्षेत्र सँकरे होते हैं और
Z चकतियाँ पास आती हैं; A पट्टी (यानी मोटे तंतु) अपनी लंबाई रखती है।
\end{solution}

\begin{solution}{exo:b2:muscle-movement:2}
(1) ATP बँधी, शीर्ष अलग; ATP का जल-अपघटन, शीर्ष तना। (2) शीर्ष
ऐक्टिन से बँधता है। (3) फ़ॉस्फ़ेट छूटा, शक्ति-आघात, ADP छूटी।
(4) नई ATP बँधती है और वह शीर्ष अलग हो जाता है। ATP उस शीर्ष को
\emph{छोड़ने} के लिए चाहिए: उसके बिना हर शीर्ष बँधा रहता है और वह
पेशी जकड़ जाती है --- यानी कठोरत्व।
\end{solution}

\begin{solution}{exo:b2:muscle-movement:3}
तंत्रिका-आवेग; ऐसिटिलकोलीन का मोचन तथा वह अंत्य-पट्ट विभव
(\qty{0.6}{ms}); उस तंतु पर तथा T नलिकाओं से नीचे पेशी का
\omterm{prop:b2:neurons-synapses:ap}{क्रिया विभव} (\qty{1}{ms}); \omterm{def:b2:cell-differentiation:fibre}{पेशीद्रव्यी जालिका} से छूटी कैल्सियम
(\qty{1}{ms}); कैल्सियम ट्रोपोनिन से बँधती है, ट्रोपोमायोसिन हटता है,
शीर्ष जुड़ते हैं (कुछ मिलीसेकंड); और बल दसियों मिलीसेकंडों भर
चढ़ता है।
\end{solution}

\begin{solution}{exo:b2:muscle-movement:4}
\omterm{def:b2:muscle-movement:motorunit}{प्रेरक इकाई}: कोई एक प्रेरक \omterm{def:b2:neurons-synapses:neuron}{न्यूरॉन} और वे सारे तंतु जिन्हें वह तंत्रिकित करता है।
\omterm{def:b2:muscle-movement:motorunit}{स्फुरण}: किसी एक आवेग से वह छोटा संकुचन। \omterm{def:b2:muscle-movement:motorunit}{समाकुंचन}: वह
जुड़ा, बना रहता संकुचन जो उन आवेगों से बनता है जो उस \omterm{def:b2:muscle-movement:motorunit}{स्फुरण} के
क्षीण होने से तेज़ पहुँचें। बल हर इकाई की दगने की दर से और
\omterm{def:b2:muscle-movement:motorunit}{भरती} हुई इकाइयों की संख्या से श्रेणीबद्ध होता है।
\end{solution}

\begin{solution}{exo:b2:muscle-movement:5}
हर अर्ध-सार्कोमियर \qty{50}{ms} में \qty{0.2}{\micro m} सरकता है:
यानी \qty{4}{\micro m/s}; यानी \qty{50}{ms} में \qty{10}{nm} के 20 आघात,
और प्रति सेकंड 400 यदि कोई एक शीर्ष पूरे समय जुड़ा रहे।
\end{solution}

\begin{solution}{exo:b2:muscle-movement:6}
\qty{0.1}{s} में \qty{2}{cm}: यानी \qty{0.2}{m/s}, प्रति सेकंड 2 लंबाइयाँ; और हर
\omterm{def:b2:cell-differentiation:fibre}{सार्कोमियर} \qty{0.1}{s} में $2/\num{40000}$ cm $= \qty{0.5}{\micro m}$ छोटा
होता है, यानी \qty{5}{\micro m/s}।
\end{solution}

\begin{solution}{exo:b2:muscle-movement:7}
$F/F_0 = 0.3125/(v/v_{\max} + 0.25) - 0.25$: प्रति सेकंड 1, 3 तथा 6
लंबाइयों पर ($0.1$, $0.3$, $0.6$ गुना $v_{\max}$): $0.64$, $0.32$, $0.12$।
शक्ति $Fv$ ($F_0\times$ लंबाइयाँ प्रति सेकंड में): $0.64$, $0.95$, $0.71$
--- यानी प्रति सेकंड 3 लंबाइयों के पास सबसे बड़ी, यानी $v_{\max}$ का लगभग एक तिहाई।
\end{solution}

\begin{solution}{exo:b2:muscle-movement:8}
$80\times 30 = \qty{2400}{N}$; पैर पर $2400\times 5/45 =
\qty{270}{N}$।
\end{solution}

\begin{solution}{exo:b2:muscle-movement:9}
\qty{10}{Hz} पर वे आवेग हर \qty{100}{ms} पर आते हैं, यानी पिछले की
कैल्सियम चले जाने के बाद: यानी अलग-अलग \omterm{def:b2:muscle-movement:motorunit}{स्फुरण}, हर एक किसी अधकचरी
ढीली दशा से शुरू होता, जो कोई लहर देता है।
\qty{50}{Hz} पर वे हर \qty{20}{ms} पर आते हैं, यानी उस कैल्सियम के
पंप किए जाने से पहले: वह कैल्सियम ऊँची बनी रहती है और वह बल जुड़ जाता है। वह
\omterm{def:b2:muscle-movement:motorunit}{समाकुंचन} उस \omterm{def:b2:muscle-movement:motorunit}{स्फुरण} से बड़ा है क्योंकि कोई एक कैल्सियम-स्पंद
अनुप्रस्थ सेतुओं के लिए श्रेणी के प्रत्यास्थ अंगों को खींचकर
पूरा बल विकसित करने भर को छोटा है; बनी रहती कैल्सियम उन्हें ऐसा करने देती है।
\end{solution}

\begin{solution}{exo:b2:muscle-movement:10}
काम \qty{20}{kJ}, ATP-ऊर्जा \qty{80}{kJ}, यानी $1.6$ मोल ATP। पहले
\qty{5}{s}: $0.8$ मोल \omterm{prop:b2:muscle-movement:energy}{क्रिएटीन फ़ॉस्फ़ेट}। बची हुई $0.8$ मोल ATP
ग्लाइकोलिसिस से: $0.4$ मोल ग्लूकोज़, $0.8$ मोल लैक्टेट --- यानी
\qty{40}{L} में, \qty{20}{mmol/L}।
\end{solution}

\begin{solution}{exo:b2:muscle-movement:11}
रुकी मोचन-वाहिका: कोई कैल्सियम नहीं, कोई संकुचन नहीं --- यानी लकवा।
रुका पंप: कैल्सियम ऊपर बनी रहती है, वह पेशी ढीली नहीं हो सकती --- यानी कोई
अनुबंधन। असंवेदी ट्रोपोनिन: ट्रोपोमायोसिन कभी हटता नहीं, कोई
संकुचन नहीं। \omterm{def:b2:blood-circulation:blood}{रक्त} में कोई कैल्सियम नहीं: कंकाल पेशी सामान्य रूप से
सिकुड़ती है (उसकी कैल्सियम भीतरी है); जबकि \omterm{def:b2:heart:anatomy}{हृदय}, जिसे अपना मोचन छेड़ने को
बाहरी कैल्सियम चाहिए, रुक जाता है।
\end{solution}

\begin{solution}{exo:b2:muscle-movement:12}
वह शीर्ष किसी ATP की मुक्त ऊर्जा का आधे तक काम में बदल देता है;
बाक़ी ऊष्मा है, और पंपों की, उन अनुप्रस्थ सेतुओं की जो गति पर बिना खींचे जुड़ते
हैं, तथा उन प्रत्यास्थ अंगों की हानियाँ उस सबको
एक चौथाई तक ले आती हैं --- यानी किसी अच्छे पेट्रोल
इंजन की दक्षता। बल श्रेणीबद्ध करने के दो रास्ते कम प्रयास पर बारीक नियंत्रण
(छोटी इकाइयाँ एक-एक करके जुड़तीं, हर एक तेज़ या
धीमे दगती) और कोई बड़ा परास (हज़ार गुना, स्फुरण-दर पर किसी एक छोटी इकाई से
\omterm{def:b2:muscle-movement:motorunit}{समाकुंचन} में हर इकाई तक), दोनों देते हैं; जबकि कोई अकेला स्विच
एक ही बल देता या कोई नहीं।
\end{solution}

\begin{solution}{pb:b2:muscle-movement:1}
\textbf{1.} $v = \sqrt{2\times 9.81\times 0.4} = \qty{2.8}{m/s}$;
$\tfrac12\times 70\times 2.8^{2} = \qty{275}{J}$।
\textbf{2.} काम $= 275 + 70\times 9.81\times 0.4 = \qty{550}{J}$,
\qty{0.4}{m} पर: यानी \qty{1370}{N}, यानी उस भार से दुगुना।
\textbf{3.} $F_0 = 500\times 30 = \qty{15000}{N}$; पैर पर
\qty{1500}{N}।
\textbf{4.} \qty{0.3}{s} में \qty{4}{cm}: यानी \qty{0.13}{m/s}, प्रति सेकंड $1.1$
लंबाइयाँ, यानी $0.14\,v_{\max}$।
\textbf{5.} $F/F_0 = 0.3125/(0.14 + 0.25) - 0.25 = 0.55$: यानी उस पेशी में \qty{8300}{N},
पैर पर \qty{830}{N} --- यानी ज़रूरी \qty{1370}{N} से कहीं
कम। केवल पेशी का छोटा होना यह छलाँग नहीं कर सकता।
बाक़ी वह प्रति-गति देती है: ज्यों ही वह छलाँगने वाला झुकता है, त्यों ही
खिंची कंडराएँ प्रत्यास्थ ऊर्जा संचित करती हैं, और वह पेशी, ऊँचे बल पर
लगभग समआयामी सिकुड़ती हुई, उन्हें \omterm{def:b2:muscle-movement:motorunit}{भरती} है; और फिर वे कंडराएँ
किसी भी तंतु के छोटे होने से तेज़ सिकुड़ती हैं।
\textbf{6.} $550/0.3 = \qty{1.8}{kW}$; यानी प्रति किलोग्राम प्रसारक \qty{92}{W/kg}।
\textbf{7.} प्रति \omterm{def:b2:cell-differentiation:fibre}{सार्कोमियर} $4/\num{48000}$ cm $= \qty{0.83}{\micro m}$,
प्रति अर्ध \qty{0.42}{\micro m}: यानी \qty{10}{nm} के 42 आघात।
\textbf{8.} \qty{0.3}{s} में \qty{0.42}{\micro m}: यानी \qty{1.4}{\micro m/s};
और लगातार जुड़े रहने पर प्रति सेकंड 140 आघात।
\textbf{9.} $500\times 6\times 10^{10}\times \num{48000}\times 300 =
4.3\times 10^{20}$ शीर्ष।
\textbf{10.} $0.55\times \num{15000}\times 0.04 = \qty{330}{J}$; प्रति
आघात $0.4\times 8\times 10^{-20} = \qty{3.2e-20}{J}$; यानी $1.0\times
10^{22}$ आघात।
\textbf{11.} $1.0\times 10^{22}/4.3\times 10^{20} = 24$ आघात प्रति
शीर्ष, उपलब्ध 42 में से: यानी लगभग \qty{60}{\%}। वह भंडार
छोटा है --- वह पेशी अपनी सीमा के पास काम कर रही है, जैसा प्रश्न 5 ने
पाया।
\textbf{12.} $1.0\times 10^{22}/6.0\times 10^{23} = 0.017$ मोल,
\qty{8.7}{g}; उसकी ऊर्जा $0.017\times 50 = \qty{860}{J}$, जिसमें से
\qty{330}{J} काम बनी: यानी \qty{39}{\%}।
\textbf{13.} \qty{30}{ms} में एक आवेग से तेज़: यानी
\qty{33}{Hz} से ऊपर; और \qty{0.3}{s} में लगभग 10 आवेग।
\textbf{14.} $9.9\times 10^{-6}\times 5\times 10^{-10}\times 6\times
10^{23} = 3\times 10^{9}$ मुक्त आयन; यानी $1.5\times 10^{9}$ ATP।
\textbf{15.} तंतु: $500/(\pi\times(3\times 10^{-3})^{2}) = 1.8\times
10^{7}$; प्रति तंतु प्रति आवेग अनुप्रस्थ-सेतु ATP $1.0\times
10^{22}/(1.8\times 10^{7}\times 10) = 6\times 10^{13}$, यानी उस पंप-आकलन से चालीस हज़ार
गुना। वह मुक्त कैल्सियम कोई छोटा शेष है: वह
जालिका कोई \qty{0.1}{mmol/L} की कोटि में छोड़ती है, यानी उस मुक्त चढ़ाई से दस
गुना, जिसका लगभग सारा तुरंत ट्रोपोनिन तथा उन बफ़रों से बँध जाता है
--- और वह सारा वापस पंप करना पड़ता है।
\textbf{16.} किसी अधिकतम छलाँग के लिए, कुछ दसियों मिलीसेकंडों के भीतर लगभग हर
इकाई, पहले छोटी और अंत में बड़ी; जबकि कोई हल्का
क़दम शायद उनका दसवें से पाँचवें भाग \omterm{def:b2:muscle-movement:motorunit}{भरती} करता है।
\textbf{17.} इकाइयाँ आकार के क्रम में \omterm{def:b2:muscle-movement:motorunit}{भरती} होती हैं: कम प्रयास पर हर
नई इकाई कोई छोटी वृद्धि जोड़ती है, और ऊँचे प्रयास पर बची हुई इकाइयाँ
बड़ी हैं और हर एक कोई बड़ा क़दम जोड़ती है।
\textbf{18.} हर तंतु का अंत्य-पट्ट विभव आधा हो जाता है; और जिन
तंतुओं का विभव देहली से नीचे गिरे वे दगते नहीं, इसलिए वह
\omterm{def:b2:muscle-movement:motorunit}{स्फुरण} तथा \omterm{def:b2:muscle-movement:motorunit}{समाकुंचन} छोटे होते हैं, और वह छलाँग छोटी पड़ती है।
\textbf{19.} $5\times 20 = \qty{0.1}{mol}$ ATP: यानी छह छलाँगों भर।
\textbf{20.} $20\times 20 = \qty{0.4}{mol}$ \omterm{prop:b2:muscle-movement:energy}{क्रिएटीन फ़ॉस्फ़ेट}:
यानी लगभग 23 छलाँगें।
\textbf{21.} दस छलाँगें $0.17$ मोल काम में लेती हैं, यानी \omterm{prop:b2:muscle-movement:energy}{क्रिएटीन
फ़ॉस्फ़ेट} अभी चुका नहीं; और कोई बीस के बाद ग्लाइकोलिसिस सँभालता है, तथा लैक्टेट
और प्रोटॉन जमा होते हैं।
\textbf{22.} \qty{40}{\%} पर \qty{1.8}{kW} काम यानी ATP का
\qty{4.6}{kW}, यानी प्रति सेकंड $0.09$ मोल, जिसे प्रति सेकंड $0.015$ मोल ऑक्सीजन चाहिए ---
यानी \qty{0.34}{L/s}, \qty{21}{L/min}, यानी \omterm{def:b2:blood-circulation:blood}{रक्त} जो ला सकता है उससे दस गुना:
वह छलाँग विवशता से अवायवीय है, और वह ऑक्सीजन बाद में आती है।
\textbf{23.} $0.17$ मोल ATP को $0.029$ मोल ऑक्सीजन चाहिए,
यानी \qty{0.64}{L}; मिनट भर में चुकाई जाए तो \qty{640}{mL/min}, यानी उस
विश्राम-खपत से ढाई गुना।
\textbf{24.} ATP \omterm{prop:b2:muscle-movement:energy}{क्रिएटीन फ़ॉस्फ़ेट} से मिलीसेकंडों के भीतर फिर
फ़ॉस्फ़ोरिलित हो जाती है, इसलिए उसकी सांद्रता मुश्किल से गिरती है; और कठोरत्व को ATP
लगभग शून्य चाहिए, जहाँ कोई जीवित पेशी किसी भी भंडार के साथ कभी नहीं पहुँचती।
\textbf{25.} प्रति अर्ध-सार्कोमियर 42 आघात उपलब्ध, 24 ज़रूरी;
उस छलाँग की गति पर $F/F_0 = 0.55$; और ATP की $0.017$ मोल, \qty{8.7}{g}।
\end{solution}
